\chapter*{导言：从对称性到可计算的物理语言}
\addcontentsline{toc}{chapter}{导言：从对称性到可计算的物理语言}
\markboth{导言}{导言}

物理学中所谓“对称性”，首先是一类\emph{操作}，其次才是一句“看起来不变”的描述。若一个操作把系统从一个状态带到另一个实验上不可区分的状态，我们就希望把这些操作放在一起研究；一旦要求它们能够复合、存在恒等操作并且每个操作可逆，这个集合便自然具有群的结构。群论的价值不在于给熟悉的几何图形重新命名，而在于把“哪些操作允许发生、这些操作怎样复合、它们怎样作用在物理态上”压缩成一套可以计算的代数语言。

对本课程而言，真正贯穿全学期的问题只有一个：
\[
\boxed{\text{系统的对称操作怎样限制可观测的物理结果？}}
\]
回答这个问题需要连续完成三次抽象。第一步把操作本身组织成群 $G$；第二步把抽象群元变成向量空间上的线性算符，即构造表示
\[
\rho:G\longrightarrow \GL(V);
\]
第三步把表示论的结论重新翻译回物理语言，例如能级简并、允许跃迁、晶体能带、角动量耦合以及全同粒子的交换对称性。本书的八章正是沿着这三步展开，而不是八个彼此独立的专题。

\chapref{chap:group-basics}与\chapref{chap:representation-theory}建立数学骨架。\chapref{chap:group-basics}回答“什么是群以及群内部怎样分解”，\chapref{chap:representation-theory}回答“群怎样在线性空间上作用以及怎样分解为不可约成分”。\chapref{chap:point-space-groups}把抽象语言第一次完整落到空间几何上：点群、晶体点群、Bravais 晶格与空间群。\chapref{chap:group-theory-quantum}再把表示论直接作用到量子力学，得到哈密顿量分块、简并、微扰劈裂、投影算符、选择定则与 Bloch 理论。\chapref{chap:rotation-group}把离散对称性推广到连续转动，对 $SO(3)$、$SU(2)$、旋量、双群和 Clebsch--Gordan 系数作系统处理。\chapref{chap:permutation-group}则研究另一类与“空间转动”本质不同的对称性：全同粒子的置换，并用 Young 图和 Young 对称化把多电子波函数的交换结构算出来。\chapref{chap:lorentz-poincare}把紧群推广到非紧群，系统地处理洛伦兹群 $SO(3,1)$、它的复化分解、旋量与 Dirac 代数、Poincar\'e 群半直积以及单粒子态的 Wigner 小群分类，从而把质量与自旋的群论来源讲清楚。\chapref{chap:lie-standard-model}再把李群、李代数的一般表示论（根、权、最高权）系统化，用于 $SU(3)$ 的夸克与色、$SU(3)_c\times SU(2)_L\times U(1)_Y$ 标准模型、电荷量子化与 Higgs 机制。四个附录分别承担晶体点群特征标、空间群速查、双群特征标重建与 Young 方法技术证明，使正文保持连续而又不牺牲完整性。\par

课程中的数学层次有意保持单向：先定义对象，再建立一般定理，再进入具体群与物理系统。一个例子只有在它能够检验或使用前面的理论时才出现；一个公式只有在后文确实需要它时才值得被强调。因此读者不应把特征标表、Schoenflies 符号或 Clebsch--Gordan 系数当成孤立记忆材料。它们分别是“共轭类上的类函数”“点群生成结构”“张量积不可约分解”这些一般理论的坐标化结果。

\section*{预备知识与阅读方式}
\addcontentsline{toc}{section}{预备知识与阅读方式}

这门课不要求读者事先学过抽象代数。我们会把“集合”“映射”“复合”“一一映射”“线性空间”等词当作真正需要使用的数学对象，而不是默认它们已经变成熟练术语。有限维线性代数仍然是必要背景，但每当基、线性算符、相似变换、本征空间、迹或张量积第一次承担新的群论含义时，正文都会重新说明它在这里为什么出现。涉及量子力学时，只需要知道态由 Hilbert 空间中的向量描述、可观测量由算符描述以及本征值方程的基本含义；涉及晶体学时，也只需要知道晶体具有周期性。倒格子、Bloch 波以及空间群所需的结构都会从定义重新建立。

阅读时最重要的是区分三个层面。若写 $g\in G$，$g$ 是抽象群元；若写 $\rho(g)$ 或 $U(g)$，它是表示空间上的线性算符；若再选定基写 $D(g)$，它才是一个具体矩阵。把这三层混为一谈，是初学表示论最常见的符号错误之一。类似地，空间转动 $R\in SO(3)$、量子态上的转动算符 $U(R)$ 与某个不可约表示矩阵 $D^{(j)}(R)$ 也应始终区分。

后文的公式中会频繁出现诸如 $\forall$、$\exists$、$\in$、$\subseteq$、$\mapsto$、$\ker$、$\operatorname{im}$ 等符号。第一次遇到时不要把它们当作需要背诵的“群论符号”：$\forall$ 读作“对任意”，$\exists$ 读作“存在”，$x\in X$ 表示“$x$ 是集合 $X$ 的元素”，$A\subseteq B$ 表示“$A$ 的每个元素都属于 $B$”，$x\mapsto f(x)$ 描述映射把输入 $x$ 送到输出 $f(x)$。$\ker$ 与 $\operatorname{im}$ 分别表示映射的核与像，等到真正需要它们时会重新定义。数学符号的目的只是把一句较长的话压缩成不含歧义的形式；如果某个符号还没有对应到一句可以说清楚的人话，就不应继续往下算。

本讲义的证明并不追求把每个成熟的线性代数恒等式都机械展开，而把篇幅留给真正决定结构的步骤：陪集分割、商群良定义、Schur 引理、有限群完全可约性、矩阵元正交、点群分类的双重计数、Seitz 乘法、$SU(2)\to SO(3)$ 二重覆盖、角动量不可约表示、Clebsch--Gordan 分解、Young 对称化与全同粒子的交换约束。读者若能独立复算这些主线推导，本课程的数学骨架就已经建立起来。

\section*{一学期的组织}
\addcontentsline{toc}{section}{一学期的组织}

若按十六至十八周组织课程，可以把\chapref{chap:group-basics}与\chapref{chap:representation-theory}安排在前五周，\chapref{chap:point-space-groups}约三周，\chapref{chap:group-theory-quantum}约三周，\chapref{chap:rotation-group}与\chapref{chap:permutation-group}各约两周，\chapref{chap:lorentz-poincare}约两周，\chapref{chap:lie-standard-model}约两周，并把最后一周用于综合例题与特征标表、空间群、双群的查表训练。这样的时间分配不是按章节页数平均，而是按认知难度分配：\chapref{chap:group-basics}和\chapref{chap:representation-theory}必须建立牢固语言，\chapref{chap:point-space-groups}与\chapref{chap:group-theory-quantum}则把它反复用于物理；\chapref{chap:rotation-group}与\chapref{chap:permutation-group}虽然数学更深，但许多结构已经能复用前面的表示论；\chapref{chap:lorentz-poincare}与\chapref{chap:lie-standard-model}把表示论从紧群推广到非紧群与一般 Lie 群，直接抵达粒子物理标准模型的群论骨架。

学完整本书后，读者应当能够从一个新问题出发，先辨认对称群，再确定相关表示空间，利用共轭类与特征标找到不可约分解，最后把分解结果翻译为物理结论。能够完成这条链条，比记住若干特征标表本身重要得多。
